% ============================================================
% Introduction section -- meant to be \input{} or \include{}'d from a
% main manuscript file, not compiled standalone. Your main.tex needs:
%   \bibliography{references}          (or \addbibresource for biblatex)
%   \bibliographystyle{<journal style>}
% pointing at ../EOACSO/references.bib, and a bibliography backend
% (natbib or biblatex) matching your target journal's citation style --
% swap \cite{} below for \citep{}/\parencite{} etc. if your class needs it.
%
% CITATION STATUS:
%   All \cite{...} calls below point at real, verified entries in
%   references.bib (see that file for source/notes on each -- including
%   one caught citation error in a source paper's own reference list).
%   Additions folded in on top of the original pass, all independently
%   verified against the publisher page / PubMed / Crossref (not taken
%   from memory):
%     - Luo:2025 (GBD 2021 update, concrete epidemiology numbers)
%     - Amato:2023 (comprehensive ML/voice-PD survey, anchors the
%       "substantial body of work" sentence)
%     - Rusz:2024 (Neurosci. Biobehav. Rev. review specifically on
%       prodromal/early digital speech biomarkers -- reinforces the
%       "vocal changes can precede motor signs" claim with a recent,
%       authoritative secondary source rather than a single primary study)
%     - BarreraGarcia:2023 (general systematic review of metaheuristic
%       feature selection, broadens the "increasingly popular" claim
%       beyond just the 5 PD-specific baseline papers already cited)
%     - Ganesan:2025 (Nov. 2025, Journal of Voice -- PD voice screening
%       via beluga whale optimization feature selection; volume/pages
%       not yet assigned as of the source check, article is online-first)
%     - Gunduz:2025 (Dec. 2025, PeerJ Comput. Sci. -- a concrete, very
%       recent illustration of the joint feature+classifier gap: pairs
%       feature selection with classifier/ensemble optimization, but as
%       two separate sequential stages rather than one joint search,
%       which is exactly the distinction this paper's contribution
%       rests on)
%   You should still independently verify each of these against the
%   original paper before submission, as with any citation.
% ============================================================

\section{Introduction}
\label{sec:introduction}

Parkinson's disease (PD) is the second most common neurodegenerative
disorder worldwide, resulting from progressive degeneration of
dopaminergic neurons in the substantia nigra \cite{Kalia:2015}. Its global
burden has risen sharply alongside population aging: an estimated 3.15
million people were living with PD in 1990, a figure that had grown to
11.8 million by 2021 and is projected to reach 25.2 million by 2050
\cite{GBD2016PDCollaborators:2018,Luo:2025}. Diagnosis remains primarily
clinical, based on the observation of cardinal motor signs ---
bradykinesia, rigidity, resting tremor, and postural instability ---
under established frameworks such as the Movement Disorder Society (MDS)
Clinical Diagnostic Criteria \cite{Postuma:2015}, typically supported by
the patient's response to levodopa therapy. This process is inherently
subjective and is particularly susceptible to misdiagnosis in early
disease stages, when motor signs may be mild, intermittent, or absent
altogether \cite{Rizzo:2016}. Because pharmacological and neuroprotective
interventions are most effective when initiated before extensive
neuronal loss has occurred \cite{Frequin:2024}, there is a pressing
clinical need for objective, low-cost, and non-invasive markers that can
support earlier and more consistent detection than motor examination
alone.

Among candidate biomarkers, vocal impairment has attracted particular
attention. Hypokinetic dysarthria is estimated to affect the large
majority of individuals with PD \cite{Ho:1999,Cao:2025}, arising from
rigidity and reduced range of motion in the laryngeal, respiratory, and
articulatory musculature. These physiological changes produce measurable
perturbations in the acoustic signal: increased cycle-to-cycle variation
in fundamental frequency (jitter), increased cycle-to-cycle variation in
amplitude (shimmer), a reduced harmonics-to-noise ratio (HNR) reflecting
incomplete glottal closure, and altered nonlinear measures of vocal-fold
vibration such as recurrence period density entropy and detrended
fluctuation analysis \cite{Yang:2020}. Such vocal changes can emerge early
in the disease course, at times preceding or co-occurring with the
earliest overt motor signs by as much as a decade
\cite{Cao:2025,Rusz:2024}, and --- unlike gait or tremor assessment ---
can be captured remotely, repeatedly, and inexpensively with nothing more
than a standard microphone. This makes voice analysis a particularly
attractive candidate for scalable, low-burden screening, including
telemedicine and home-monitoring scenarios that cannot easily
accommodate in-person motor examination \cite{Rusz:2024}.

Motivated by this, a substantial and rapidly growing body of work has
applied machine learning to acoustic features extracted from
sustained-phonation or connected-speech recordings for automated PD
detection \cite{Amato:2023}, with reported accuracies frequently exceeding
90\% on public benchmark datasets such as the Oxford voice dataset
introduced by Little et al.~\cite{Little:2008}. Because the number of
extractable acoustic measures is often large relative to the number of
available recordings, feature selection is commonly used to reduce
dimensionality, limit overfitting, and improve interpretability. Wrapper
approaches built around population-based metaheuristic optimization have
become an increasingly popular choice for this task, both across
biomedical feature selection generally \cite{BarreraGarcia:2023} and
specifically for PD voice screening: particle-swarm-, butterfly-, and
whale-inspired variants \cite{Hashemi:2026}, a whale-to-grey-wolf cascade
\cite{AlNajjar:2024}, a modified grey wolf optimizer \cite{Santhosh:2025},
a modified hunger games search \cite{Hashim:2023}, a quantum mayfly
optimizer \cite{Mansour:2024}, and a beluga whale optimizer
\cite{Ganesan:2025}.
Yet the large majority of this literature treats the choice of classifier
as fixed and pre-determined, optimizing only the feature subset; the
complementary question of which classifier (or combination of
classifiers) best exploits a given feature subset is rarely optimized
jointly within the same search process. Even recent work that pairs
feature selection with classifier-side optimization typically does so as
two separate, sequential stages --- e.g., a filter-based feature-selection
step followed by a separately optimized classification ensemble --- rather
than as a single joint search over both design choices at once
\cite{Gunduz:2025}.

To address this limitation, this study proposes a joint
feature-and-classifier optimization pipeline for voice-based PD
screening, built on an Elite-guided Opposition-based Archive Competitive
Swarm Optimizer (EOACSO). The main contributions of this work are as
follows:

\begin{itemize}
  \item We formulate PD voice screening as a joint search over the
    acoustic feature subset and the classifier used to exploit it, within
    a single optimization encoding, rather than optimizing feature
    selection alone.
  \item We evaluate every method under a subject-grouped cross-validation
    protocol that prevents recording-level data leakage, yielding
    performance estimates that better reflect realistic,
    subject-independent deployment.
  \item We re-implement eight recently published metaheuristic
    feature-selection methods spanning five source papers --
    particle-swarm/butterfly/whale-based variants \cite{Hashemi:2026}, a
    whale-to-grey-wolf cascade \cite{AlNajjar:2024}, a modified grey wolf
    optimizer \cite{Santhosh:2025}, a modified hunger games search
    \cite{Hashim:2023}, and a quantum mayfly optimizer \cite{Mansour:2024}
    -- and evaluate them, together with the proposed method, under one
    shared encoding, fitness function, and evaluation protocol, enabling
    a controlled, like-for-like comparison that is largely absent from
    the current literature.
  \item We conduct a systematic ablation study isolating the individual
    contribution of each of EOACSO's three constituent strategies:
    elite-guided update, opposition-based reinitialization, and a
    diversity-preserving elite archive.
  \item We examine the acoustic features most consistently selected by
    the proposed pipeline in light of established physiological
    correlates of parkinsonian dysphonia, to assess the biological
    plausibility --- and potential clinical interpretability --- of the
    resulting screening model.
\end{itemize}

The remainder of this paper is organized as follows.
Section~\ref{sec:2} reviews related work on PD diagnosis,
voice biomarkers, and metaheuristic feature selection.
Section~\ref{sec:3} describes the datasets and the proposed
methodology. Section~\ref{sec:results} presents the experimental
results, including comparative performance, statistical significance
testing, and the ablation study. Section~\ref{sec:discussion} discusses
the clinical implications and limitations of the proposed approach.
Section~\ref{sec:conclusion} concludes the paper.
